\chapter{Apr.~22 --- Geometric Aspects, Part 2}

\section{Relation to KZ Equations}

\begin{remark}
  Recall that for
  $V = M_{\mu_1} \otimes \cdots \otimes M_{\mu_N}$
  and $z_1, \dots, z_N$, we have
  \[
    \psi_C
    = \sum_{\vec{m}}
    \bigg(\int_C \psi_m(\vec{z}, \vec{t})
    \rho_m(\vec{z}, \vec{t})\, d\vec{t}\bigg)
    v_{\vec{m}},
  \]
  where $\vec{m} = (m_1, \dots, m_N)$
  with $\sum m_i = m$ and
  $v_{\vec{m}} = e^{m_1} v_1 \otimes \cdots \otimes e^{m_N} v_N$.
  Note that $\dim Y_{\vec{z}, m} = m$.
  Also, the above integral depends only
  on the class of
  $C \in C_m(Y_{\vec{z}, m}, \mathcal{L}^*)$.
\end{remark}

\begin{theorem}[Aomoto]
  Let $m = 1$. Then for almost
  all $\mu, \kappa$,
  \begin{enumerate}
    \item $H_0(Y_{\vec{z}, 1}, \mathcal{L}^*) = 0$;
    \item $\dim H_1(Y_{\vec{z}, 1}, \mathcal{L}^*) = \dim H^1(Y_{\vec{z}, 1}, \mathcal{L}) = N - 1$;
    \item the forms
    $\psi(\vec{z}, t) dt / (t - z_j)$
    for $j = 1, \dots, N$
    for a basis in
    $H^1(Y_{\vec{z}, 1}, \mathcal{L})$;
  \item the Pochhammer
    loops $C_i$
    around $z_i, z_N$ form a basis
    in $H_1(Y_{\vec{z}, 1}, \mathcal{L}^*)$.
  \end{enumerate}
\end{theorem}

\begin{corollary}\label{cor:level-1}
  For every $\vec{z}$ and almost all
  $\mu_i, \kappa$, the
  map $C \mapsto \psi_C$ is an isomorphism
  between $H_1(Y_{\vec{z}, 1}, \mathcal{L}^*)$
  and the space of $W$-valued solutions
  of the KZ equations, where
  $W$ is the space of singular
  vectors of weight $-\sum \mu_i + 2$.
\end{corollary}

\begin{theorem}[Arbitrary level]\label{thm:arbitrary-level}
  For all $\vec{z}$ and
  almost all $\vec{\mu}, \kappa$,
  \begin{enumerate}
    \item $H_i(Y_{\vec{z}, m}, \mathcal{L}^*) = 0$ for $i < m$;
    \item $\dim H_m(Y_{\vec{z}, m}, \mathcal{L}^*) = \dim H^m(Y_{\vec{z}, m}, \mathcal{L}) = m! p_{N - 1}(m)$,
      where $p_k(m)$ is the number of partitions
      of $m$ into ordered sums of
      $k$ nonnegative integers.
  \end{enumerate}
\end{theorem}

\begin{remark}
  Taking the smaller space
  $H_m(Y_{\vec{z}, m}, \mathcal{L}^*)^{S_m}$,
  gets rid of the $m!$
  in Theorem \ref{thm:arbitrary-level}(2).
  This gives a natural generalization of
  Corollary \ref{cor:level-1} for arbitrary
  $m$: We have
  \[
    H_m(Y_{\vec{z}, m}, \mathcal{L}^*)^{S_m}
    \cong W,
  \]
  where $W$ is the space of singular
  vectors of weight
  $-\sum m_i + 2m$.
\end{remark}

\section{Gauss-Manin Connection}

\begin{remark}
  Consider the fiber bundle 
  $\pi : X_{m + N} \to X_N$
  with fiber $Y_{\vec{z}, m}$.
  Recall that there is an action of $S_N$
  on $X_N$. We have can choose a cycle in
  every fiber, but we want to match them in
  a compatible way.

  Note that
  $\psi_m(\vec{z}, \vec{t})$
  defines a local system
  on $X_{m + N}$, which gives local
  systems on $Y_{\vec{z}, m}$
  by restriction.
  Thus we get a vector bundle
  $H_m(X_{m + N} / X_N, \mathcal{L})$
  with a fiber at $\vec{z}$
  given by $H_m(Y_{\vec{z}, m}, \mathcal{L})$.
  Note that by analytic continuation we may
  move $\vec{z} \to \vec{z}'$, so for any
  $\vec{z} \in X_N$,
  a choice of
  $C \in H_i(Y_{\vec{z}, m}, \mathcal{L})$
  defines a cycle in each
  $Y_{\vec{z}', m}$ defines a cycle in
  each $Y_{\vec{z}', m}$ for $z'$
  close enough to $z$.
  This defines a connection
  in
  \[
    C(X_{m + N} / X_N, \mathcal{L})
  \]
  via the condition that all local sections
  produced in this way are flat
  (the \emph{Gauss-Manin connection}).
\end{remark}

\begin{theorem}\label{thm:gauss-manin-kz}
  For all $\mu_i$ and almost all $\kappa$,
  the map $C \mapsto \psi_C$ gives an
  isomorphism of local systems
  \[
    H_m(X_{m + N} / X_N, \mathcal{L}(\mu, \kappa))^{S_m}
    \cong W_{\mathrm{KZ}},
  \]
  where the left-hand side is endowed by
  the Gauss-Manin connection
  and the right-hand side is the
  trivial vector bundle over
  $X_N$ with fiber
  $W = (V^{\mathfrak{n}^-})^\lambda$
  ($\lambda = -\sum \mu_i + 2m$),
  endowed with the KZ connection.
\end{theorem}

\begin{remark}
  Theorem \ref{thm:gauss-manin-kz}
  fails for $k \in \Z_+$
  (recall that $\kappa = k + h^\vee$).
\end{remark}

\section{Relative Homology}

\begin{example}\label{ex:pochhammer}
  Let $X = \C \setminus \{0, 1\}$ and
  $E$ the $1$-dimensional local system
  on $X$ defined by
  \[
    \psi(z) = z^\alpha(1 - z)^\beta
  \]
  for $\alpha, \beta \in \C$. Then for
  generic $\alpha, \beta$, we have
  $H_1(X, E^*) \cong \C$, generated
  by the Pochhammer loop.
  For $\re \alpha \gg 0$
  and $\re \beta \gg 0$, we can write
  \[
    \int_C f(z) \psi(z)\, dz
    = (1 - e^{2\pi i \alpha})(1 - e^{2\pi i \beta})
    \int_0^1 f(z)\, \psi(z)\, dz.
  \]
  We call this integral on the right-hand side
  the \emph{regularization} of $[0, 1]$,
  and denote it
  by $[0, 1]_{\mathrm{reg}}$.
\end{example}

\begin{definition}
  Let $X = \overline{X} \setminus D$,
  where $\overline{X}$ is a real manifold
  and $D \subseteq \overline{X}$ is a
  closed subset.
  The space
  $C_k(\overline{X}, D, E)$ of
  \emph{relative chains}
  is the space of finite linear combinations
  \[
    \sum \varepsilon_i \sigma_i,
  \]
  where $\varepsilon_i$ is a
  flat section of $E$ over
  $\sigma_i(\Sigma^k) \cap X$ and
  $\sigma_i$ is a singular $k$-simplex
  in $\overline{X}$.
\end{definition}

\begin{remark}
  Note that
  $C_k(X, E) \subseteq C_k(\overline{X}, D, E)$,
  which gives a map
  \[
    H_k(X, E)
    \longrightarrow H_k(\overline{X}, D, E),
    \tag{$*$}
  \]
  though this map is
  generally not injective or surjective.
\end{remark}

\begin{example}
  Note that
  $[C] = (1 - e^{2\pi i \alpha})(1 - e^{2\pi i \beta}) [0, 1]$
  in Example \ref{ex:pochhammer}, and
  in fact the map on homology is surjective
  for all $\alpha, \beta \in \Z$.
  This is not true for standard homology,
  e.g. consider $\alpha = \beta = 0$.
\end{example}

\begin{theorem}
  For a general hyperplane arrangement
  and general $\lambda_H$, the
  map $(*)$
  is an isomorphism.
\end{theorem}

\begin{definition}
  Call any subset $L$ of the set of
  all hyperplanes $S$ such that
  $\bigcap_{H \in L} H \ne \varnothing$
  an \emph{edge}, and denote
  $\lambda_L = \sum_{H \in L} \lambda_H$.
\end{definition}

\begin{theorem}
  If for every edge $L$,
  we have $\lambda_L \notin \Z$,
  then the map $(*)$ is an isomorphism.
\end{theorem}

\begin{theorem}
  Let $D_{\vec{z}, m}$ be the
  union of the hyperplanes
  $\{t_i = t_j, t_i = z_k\}$, so that
  $Y_{\vec{z}, m} = \C^m \setminus D_{\vec{z}, m}$.
  Then for all $\mu_i, \kappa$ generic,
  we have:
  \begin{enumerate}
    \item $H_i(\C^m, D_{\vec{z}, m}, \mathcal{L}(\vec{\mu}, \kappa)) = 0$
      for $i < m$.
    \item For $\vec{z} \in \R^N$,
      define
      \[
        Y_{\vec{z}, m}^\R
        = \{\vec{t} \in \R^m : t_i \ne t_j, t_i \ne z_k\}
      \]
      Let $\{D_i\}$ be the set of bounded
      connected components of
      $Y_{\vec{z}, m}^\R$ (so each $D_i$ is
      the interior
      of some convex polytope), and fix
      an arbitrary section of
      $\mathcal{L}(\vec{\mu}, \kappa)$
      over each $D_i$. Then
      \[
        H_m(\C^m, D_{\vec{z}, m}, \mathcal{L}(\vec{\mu}, \kappa))
        = \bigoplus \C [D_i],
      \]
      and thus
      $H_i(Y_{\vec{z}, m}, \mathcal{L}(\vec{\mu}, \kappa)) = 0$
      for $i < m$.
  \end{enumerate}
\end{theorem}

\begin{remark}
  We can pick another basis for
  the homology
  from
  $\vec{m} = (m_1, \dots, m_{N - 1})$
  (with $m_i \ge 0$) and a bijection $\phi$ which assigns to a
  pair $1 \le a, b \le N - 1$ an
  index $i = \phi(a, b) \in \{1, \dots, m\}$
  ($m = \sum m_i$).
\end{remark}

\begin{theorem}
  Let $z_1, \dots, z_N$ such that
  $[z_i, z_N]$ do not intersect
  except for at $z_N$, and define
  \[
    D_{\vec{m}, \phi}
    =
    \left\{
      t_1, \dots, t_m
      : \substack{\displaystyle t_i = z_N + u^a_b(z_a - z_N)\text{ if } i = \phi(a, b),\\
      \displaystyle 0 < u_1^1 < u_2^1 < \cdots < u^1_{m_1} < 1, \\ \displaystyle \vdots \\ \displaystyle 0 < u_1^{N - 1} < \cdots < u^{N - 1}_{m_{N - 1}} < 1}
    \right\}
    \subseteq Y_{\vec{z}, m}.
  \]
  Then for almost all $\mu_i, \kappa$,
  the $D_{\vec{m}, \phi}$ form a basis
  for the relative homology.
\end{theorem}

\section{Monodromy of KZ Equations}

\begin{remark}
  Recall the KZ equations:
  \[
    \kappa \frac{\partial}{\partial z_i}
    \psi
    = \bigg(
      \sum_{j = 1, j \ne i}^N \frac{\Omega_{i, j}}{z_i - z_j}
    \bigg) \psi, \quad \kappa \notin \Q.
  \]
  We can interpret this system as a flat
  connection in a trivial bundle
  with fiber
  $V_1 \otimes \cdots \otimes V_N$ over
  \[
    X_N = \{
      (z_1, \dots, z_N) \in \C^N : z_i \ne z_j
    \}.
  \]
  The solutions are then given by
  the flat sections, which
  we denote by $\Gamma_f(U, V_{\mathrm{KZ}})$.
\end{remark}

\begin{prop}
  For every $\gamma$, the map
  $M_\gamma : V \to V$ is a $\g$-homomorphism
  and $M_\gamma$ preserves singular
  vectors. If
  $\vec{z}_0 = (z_1, \dots, z_N)$ and
  $\gamma(0) = \gamma(1) = \vec{z}_0$, then
  $\gamma \to M_\gamma$ is a
  representation of
  $\pi_1(X_N, \vec{z}_0)$.
\end{prop}

\begin{theorem}[Artin]\label{thm:artin}
  We have the following:
  \begin{enumerate}
    \item $\pi_1(X_N) = PB_N$, the
      \emph{pure braid group}
      defined by
      $PB_N = \ker \sigma$ where
      $\sigma : B_N \to S_N$;
    \item $S_N$ acts on $X_N$ by
      permuting coordinates, and
      $\pi_1(X_N / S_N) = B_N$.
  \end{enumerate}
\end{theorem}

\begin{remark}
  The map
  $B_N \to \pi_{1}(X_N / S_N)$ in
  Theorem \ref{thm:artin}(2) is defined as
  follows: Fix
  $\vec{z}_0 = (z_1, \dots, z_N) \in \R^N$
  such that $z_N < z_{N - 1} < \cdots < z_2 < z_1$,
  and for each $\sigma \in S_N$ let
  \[
    V^\sigma
    = V_{\sigma^{-1}(1)} \otimes \cdots \otimes V_{\sigma^{-1}(N)}.
  \]
  \pagebreak
  Let $\sigma : V \to V^\sigma$
  be defined by
  $v_1 \otimes \cdots \otimes v_N \mapsto
  v_{\sigma^{-1}(1)} \otimes \cdots \otimes v_{\sigma^{-1}(N)}$.
  Fix a basepoint and lift a loop
  $\gamma$ in $X_N / S_N$ to a path
  in $X_N$. We have an operator
  \[
    \widetilde{M}_\gamma
    = \sigma M_\gamma
    : V \longrightarrow V^\sigma,
  \]
  and we can define
  $\widecheck{M}^{\pm}_i(z) = (\widetilde{M}_{\gamma_i})^{\pm 1}$
  for some path $\gamma_i$.
\end{remark}

\begin{theorem}
  We have the following:
  \begin{enumerate}
    \item $\widecheck{M}^{\pm}_i \widecheck{M}^{\mp}_i = \id$;
    \item $\widecheck{M}^{\pm}_i \widecheck{M}^{\pm}_{i + 1} \widecheck{M}^{\pm}_i = \widecheck{M}^{\pm}_{i + 1} \widecheck{M}^{\pm}_i \widecheck{M}^{\pm}_{i + 1}$.
  \end{enumerate}
\end{theorem}

\begin{corollary}
  If all representations
  $V_i \cong W$ are isomorphic
  and $z_1 > \cdots > z_N$, then
  $b_i \mapsto \widecheck{M}^{\pm}_i(\vec{z})$
  is a representation of the braid
  group $B_N$ on $V = W^{\otimes N}$.
\end{corollary}
